\chapter{Elasticity}

\section{Nonlinear Elastodynamics}

The bulk modulus $B$, a measure of local compressibility of a substance, is
defined as the ratio of infinitesimal change in pressure $dP$ caused by an
infinitesimal change in volume $dV$, at a given volume $V$
\be 
B \defe - V \dst\frac{dP}{dV}.
\ee
We require $B > 0$, and since $V > 0$, then $dP/dV < 0$.
This results in behaviors observed physically, namely that, all else being
equal, 
\begin{itemize}
\item increase in pressure causes decrease in volume, and vice versa, 
$\uparrow P \Longleftrightarrow \;\downarrow V$, and
\item increase in volume causes decrease in pressure, and vice versa, 
$\uparrow V \Longleftrightarrow \;\downarrow P$.
\end{itemize}
Using the definition of density $\rho \defe m / V$, 
\be
\frac{d\rho}{dV} = \frac{-m}{V^2},
\ee
and
\be
\frac{dP}{dV} = \frac{dP}{d\rho}\frac{d\rho}{dV},
\ee
the bulk modulus $B$ may also be written as
\be
B = \rho \frac{dP}{d\rho}.
\ee
Since $B > 0$ and $\rho > 0$, then $dP/d\rho > 0$.  This result physically means
that, all else being equal, 
\begin{itemize}
\item increases in pressure cause increases in density (and vice versa), and
\item increases in density cause increases in pressure (and vice versa).
\end{itemize}

\section{Linear Elastodynamics}
Linear elastodynamics begins with nonlinear elastodynamics, and adds progressively
simplifying (and thus, restrictive) assumptions.  The first assumption will be regarding
linearity of displacement gradients.  
\begin{itemize}
\item {\bf Assumption 1:} Second-order and higher displacement gradients are 
neglible,  
\end{itemize}
which restricts the stress-response function to be a special 
case of the fully general stress response function
\be
 \tPbold = \tPbold (\tFbold).
\ee
Rewriting the foregoing stress-response function in terms of the displacement 
gradient,
\be
 \tPbold(\tFbold) = \tPbold(\vid + \Grad \vu),
\ee
gives an expression that can be cast in terms of a Taylor series expansion, 
\begin{align}
 f(x_0 + \Delta x) &= f(x_0) + f'(x_0) \Delta x + \frac{f''(x_0)}{2!} (\Delta x)^2 + 
 \mbox{\footnotesize H.O.T.}, \\
 &= \sum_{j=0}^{\infty} \frac{f^j(x_0)}{j!} (\Delta x)^j.
\end{align}
Then 
\be
 \tPbold(\vid + \Grad \vu) = \tPbold(\vid) + \frac{\partial \tPbold(\vid)}{\partial \tFbold}
 \Grad \vu + \mathcal{O}\left[(\Grad \vu)^2\right].
\ee
Next, we make another assumption,
\begin{itemize}
\item {\bf Assumption 2:} The initial state is stress free, thus $\tPbold(\vid) = \vzero$.
\end{itemize}
The fourth-order elasticity tensor $\fotC$ is defined as,
\be
 \fotC \defe \frac{\partial \tPbold(\vid)}{\partial \tFbold}.
\ee
Symmetry of the elasticity tensor provides
\be
 \fotC \; : \; \Grad \vu \;\;=\;\; \fotC \; : \; \mbox{\sym}(\Grad \vu).
\ee
Finally, 
\be 
 \mbox{\sym}(\Grad \vu) = \half(\Grad \vu + \Grad \vu\transpose) = \ve{\epsilon}.
\ee
These results give
\be
\tPbold(\vid + \Grad \vu) = \vzero \;\;+\;\; \fotC \; : \; \ve{\epsilon} \;\;+\;\; 
 \mbox{\footnotesize H.O.T.},
\ee
or simply
\be
\tPbold(\vid + \Grad \vu) \overset{{\mbox{\tiny(LIN)}}}{=} 
\tPbold_{\tiny \mbox{LIN}} (\ve{\epsilon}) = \fotC \; : \; \ve{\epsilon}.
\ee
Next, we can make another limiting assumption, 
\begin{itemize}
\item {\bf Assumption 3:} The material is {\bf isotropic},
\end{itemize}
which allows the stress-repsonse function to be written in terms 
of two Lam\'{e} parameters $\lambda$ and $\mu$ as
\be
 \tPbold_{\mbox{\tiny{LIN}}} (\ve{\epsilon}) 
 = \lambda \; \trace(\ve{\epsilon}) \; \vid + 2 \; \mu \; \ve{\epsilon}.
\ee
Next, we make another assumption, 
\begin{itemize}
\item {\bf Assumption 4:} The material is {\bf homogeneous},
\end{itemize}
thus $\rho_0$, $\lambda$, and $\mu$ are constants with respect to the position $\vX$ in the
body.

In the conservation of momentum equation, the divergence of the stress tensor appears.
Thus, we find
\begin{align}
\Grad \cdot \tPbold_{\mbox{\tiny{LIN}}} (\ve{\epsilon})
&= \Grad \cdot \left(
\lambda \; \trace(\ve{\epsilon}) \; \vid \;\;+\;\; 2 \; \mu \; \ve{\epsilon}
\right) \\
&= \lambda \; \Grad \left(\Grad \cdot \vu \right) \;\;+\;\; 
\mu \left[ 
 \Grad \cdot \left(\Grad \vu \right) + \Grad \left(\Grad \cdot \vu \right)
\right], \\
&= (\lambda + \mu) \; \Grad \left(\Grad \cdot \vu \right) \;\;+\;\; \mu \; \Grad \cdot \left(\Grad \vu \right).
\end{align}
Now, since
\be
\Grad \cdot \left( \Grad \vu \right) = \Grad \left( \Grad \cdot \vu \right) - 
\left( \Grad \times \Grad \times \vu \right),
\ee
we can write
\be
\Grad \cdot \tPbold_{\mbox{\tiny{LIN}}} (\ve{\epsilon})
 = \underset{\mbox{\footnotesize dilitational}}{\underbrace{(\lambda + 2\mu) \; \Grad \left(\Grad \cdot \vu \right)}}
 \;\;-\;\; 
 \underset{\mbox{\footnotesize rotational}}{\underbrace{\mu \; \left(\Grad \times \Grad \times \vu \right)}}
\ee
Finally, we make another assumption
\begin{itemize}
\item {\bf Assumption 5:} The {\bf body force} is 
{\bf negligible}, thus $\vb_0 = \vzero$.
\end{itemize}
This final assumption allows the {\bf progressive wave equation for solids} to be 
written as
\be
 \rho_0 \; \ddot{\vu}  
 = 
 \underset{\mbox{\footnotesize dilitational}}{\underbrace{(\lambda + 2\mu) \; \Grad \left(\Grad \cdot \vu \right)}}
 \;\;-\;\; 
 \underset{\mbox{\footnotesize rotational}}{\underbrace{\mu \; \left(\Grad \times \Grad \times \vu \right)}}
\ee
There are two wave speeds: the longitudinal sound speed $c_L$ and the transverse
sound speed $c_T$.  
\begin{align}
 c_L &= \sqrt{\dst\frac{\lambda + 2 \; \mu}{\rho_0}} \hspace{1.0cm} \mbox{\bf (longitudinal
 or P-wave)}, \\
 c_T &= \sqrt{\dst\frac{\mu}{\rho_0}} \hspace{1.9cm} \mbox{\bf (transverse of S-wave)}.
\end{align}
\begin{Hremark}
The longitudinal sound speed is faster than the transverse sound speed.
\end{Hremark}

The bulk modoulus $B$ is related to the Lam\'{e} constants $\lambda$ and $\mu$ through
\be
B = \lambda + \frac{2\mu}{3}.
\ee
Substituting for $B$ in the longitudinal sound speed gives
\be
c_L = \sqrt{\dst\frac{B + 
\frac{4}{3} \; \mu }{\rho_0}} \hspace{1.0cm} \mbox{\bf (longitudinal or P-wave)}.
\ee
\begin{Hremark}
Often, the Lam\'{e} constant $\mu$ will be denoted $G$, which is the 
shear modulus.  Preference for $\mu$ or $G$ is a matter of taste.  
\end{Hremark}














